\subsection{与本项目相关的研究工作积累}
申请人长期从事时间序列分析、新能源功率预测及多尺度特征表示研究，围绕光伏功率的日周期与局部波动，积累了小波分解、非线性映射和时间依赖建模经验。本项目拟进一步研究气象信息获取时限下的概率预测与爬坡风险；专属先导实验尚无可核验证据，区间可靠性与新疆场站预警效果需独立检验。
% F-XJ-07：尚未提供本项目专属先导数据、代码、日志和性能证据。

\subsection{已取得的研究工作成绩}
申请人作为通讯作者发表 DWT-Former 研究（\textit{Energy}, 2025, 139283；DOI：\href{https://doi.org/10.1016/j.energy.2025.139283}{\nolinkurl{10.1016/j.energy.2025.139283}}）\cite{sun2025dwt}，围绕小波多尺度特征与 Transformer 时间表示的融合开展光伏功率预测研究，为本项目的尺度对照提供方法基础。

申请人作为共同第一作者发表 WaveKAN 研究（\textit{Neurocomputing}, 2026, 133361；DOI：\href{https://doi.org/10.1016/j.neucom.2026.133361}{\nolinkurl{10.1016/j.neucom.2026.133361}}）\cite{liu2026wavekan}，研究自适应小波分解与 Kolmogorov--Arnold 网络结合的长时序预测方法。两项论文的具体项目归属依据成果记录核实。

\subsection{已具备的实验条件}
申请人依托新疆大学开展研究，可稳定使用 NVIDIA A40 8块和 NVIDIA GeForce RTX 3090 Ti 6块，具备 Linux、Python、深度学习框架、Git 和 LaTeX 等软件工具。这些条件能够支持数据质量检查、多尺度建模、概率预测训练和滚动评价。

申请人负责总体方案和实验结果核查，\XJField{student-count}名新疆大学研究生按成员表分担数据、建模、校准、爬坡评价与原型任务。各成员2027--2028年的实际投入时间及本人确认尚需落实。学院公共算力拟作为补充条件申请，现阶段不计入固定资源配额。
% F-XJ-03：成员投入时间、个人信息、本人签名及其他在研任务仍需核对。

\subsection{尚缺少的实验条件及拟解决的途径}
\subsubsection{新疆光伏场站实测数据}
新疆场站实测数据及使用授权尚待核实落实，拟按1个主研究场站、连续24个月历史记录的目标获取数据，连续12个月为方法验证的最低覆盖要求，时间分辨率应支持15分钟功率与气象序列。申请人在项目首季度对接数据权利人，核验提供单位、实际覆盖时段、授权范围及缺测、停机和限电标记，完成数据清单与质量评估。前期方法验证使用许可允许的公开光伏与气象数据，保留来源、版本及变量定义，其结果与新疆实测验证分别报告。目前尚无可核验的项目专属数据库材料。
% F-XJ-01：数据提供单位、实际覆盖时段及授权范围待落实。

\subsubsection{可追溯的气象预报输入}
可追溯业务 NWP 档案尚待落实，需核实来源、发布时刻、预报步长和覆盖范围，并按实际获取时刻组织回溯实验。再分析资料用于背景描述和机制探索；缺少业务预报档案时，相关结论限定为离线方法研究。
% F-XJ-02：业务NWP档案与获取条件待落实。

\subsubsection{合作验证与人员安排}
拟与气象业务单位或场站合作，获取需求反馈并核验天气个例。具体合作方、承担事项、书面协议及业务研究基地接入证明尚待落实；协议需明确数据用途、验证任务和成果使用范围。
% F-XJ-04：合作单位、承担事项与书面落实情况待确认。

申请经费3万元，其中业务费1.50万元、劳务费1.20万元、管理费0.30万元，使用现有设备开展模型试验。劳务按24人月、每人月500元测算，支付给实际承担数据整理、实验和软件任务的研究生，按工作记录及依托单位管理要求执行。

\subsection{相关在研项目}
申请人主持自治区高校基本科研业务费科研项目（培育类）“融合多源数据大模型驱动的光伏发电功率预测方法研究”，批准号XJEDU2026P005，资助金额5万元（直接经费3.9万元、间接经费1.1万元），执行期2026年1月至12月，负责总体方案、模型研究与成果组织。本申请重点研究可用气象信息约束下的概率预测、天气转换与爬坡事件。两项目的任务、数据及成果边界需依据任务书逐项核对，支出与成果归属分别记录；该在研项目的限项适用情况尚待基金管理方核实。
% F-XJ-05：在研项目限项适用、截止日前验收材料状态及任务书对照仍需核实。

申请人参与国家自然科学基金地区项目“数据驱动下基于互补型生成对抗网络与神经统计模型的公路沙害智能预警研究”，批准号62466056，资助金额32万元，执行期2025年1月至2028年12月，基金委登记单位为新疆大学，本人负责数据收集整理。该项目研究公路沙害，数据质量控制经验可为本申请提供方法参考；资料和成果的使用遵循各自授权范围。申请人尚未主持完成国家自然科学基金项目。

\XJField{student-count}名研究生成员参与其他在研项目及任务重合的情况仍需逐人核对，并结合其可投入时间落实分工。

\subsection{申请人简历}
\XJField{applicant}，\XJField{sex}，\XJField{birth}出生，\XJField{degree}，\XJField{organization}\XJField{professional-title}、硕士研究生导师。主要从事时间序列分析、新能源功率预测、多尺度时频表示与多源信息融合研究。现所在院系为\XJField{department}。通讯地址：\XJField{address}；邮编：\XJField{postcode}；电子邮箱：\XJField{email}。
% 院系名称、邮编、邮箱及以下履历年月沿用共用档案的公开来源待终核状态，不等同于申请人已终核；详见docs/填写说明与待补清单.md。

\subsubsection{教育经历}
2019.09--2023.08，新疆大学，计算机科学与技术专业，博士。

2012.09--2014.06，上海海事大学，计算机技术专业，硕士。

2008.09--2012.06，海南热带海洋学院，生态学专业，学士。

\subsubsection{工作经历}
2023.12至今，新疆大学，副教授。

2023.10--2023.12，新疆大学，教学与科研工作。

2018.03--2019.08，浙江大学山东工业技术研究院，软件工程师。

2016.10--2017.12，上海中科软科技有限公司，软件工程师。

\subsubsection{与本项目直接相关的代表成果}
申请人作为通讯作者发表DWT-Former研究（\textit{Energy}, 2025, 341: 139283）\cite{sun2025dwt}，作为共同第一作者发表WaveKAN研究（\textit{Neurocomputing}, 2026, 682: 133361）\cite{liu2026wavekan}。上述成果分别支撑光伏多尺度预测和自适应小波长时序建模的方法基础；本项目专属概率校准、爬坡预警及新疆场站先导试验尚无可核验证据。

